%% ================================
%% Dokumenteinstellungen & Metadaten
%% ================================

% Autor & Abgabe
\newcommand{\Autor}{Parjanya Tapan Dalal}
\newcommand{\Matrikelnummer}{2600529}
\newcommand{\Datum}{31.01.2026}

% Titel
\newcommand{\Titel}{Entwicklung eines automatisierbaren Bewegungssystems zum Aufbau von Linearführungen}
\newcommand{\TitelEN}{Development of an Automatable Motion System for the Assembly of Linear Guides}

% Art der Arbeit
\newcommand{\Dokumentenart}{Bachelorarbeit}
\newcommand{\Pruefungsart}{Bachelorprüfung}

% Prüfer & Betreuer
\newcommand{\Erstpruefer}{Erstpruefer: Prof. Dr.-Ing. Alexander Frenkel}
\newcommand{\Zweitpruefer}{Zweitpruefer: Bernd Lüerß}
\newcommand{\Betreuer}{Industrieller Betreuer: Bernd Lüerß}

% Firma
\newcommand{\Firma}{
    in Zusammenarbeit mit:\\
    SmarAct GmbH\\
    Abteilung Technologieplattform\\
    Schütte-Lanz-Straße 9\\
    26135 Oldenburg
}

%% ================================
%% Abstract-Informationen
%% ================================

% Stichworte (Keywords)
\newcommand{\Stichworte}{
    Automatisierung, Linearführung, Automatisierung von Fertigungsprozesssen, Einschraubroboter, Vorspannung, Steuerungstechnik, Piezo, Stick-Slip-Antriebe
}

\newcommand{\StichworteEN}{
    Automation, linear guides, manufacturing process automation, screw-driving robots, preload, control engineering, piezoelectric actuators, stick-slip drives
}

% Zusammenfassung (DE)
\newcommand{\Zusammenfassung}{
Im Rahmen dieser Bachelorarbeit wurde für die SmarAct GmbH ein automatisierbares System zur Montage von Linearführungen entwickelt. Ziel war es, den bisher manuellen Montageprozess effizienter, reproduzierbarer und qualitativ stabiler zu gestalten. Dazu wurden bestehende Montageprozesse analysiert und technische Anforderungen abgeleitet.\\
Auf dieser Basis wurde ein Montagesystem konzipiert, umgesetzt und erprobt. Zentrale Systemkomponenten sind ein präzises Positioniersystem, automatisierte Schraubwerkzeuge, ein Presssystem zur Montage der Keile sowie ein Schiebesystem zur Einstellung der Vorspannung. Ergänzend wurde ein Steuerungskonzept zur Ablaufüberwachung und Validierung entwickelt. Die Lösung wurde erfolgreich getestet und optimiert, um eine stabile und qualitativ hochwertige Serienmontage zu ermöglichen.\\
Diese Abschlussarbeit unterliegt einer Geheimhaltungsvereinbarung. Detaillierte Ergebnisse sind nicht öffentlich zugänglich.
}

% Abstract (EN)
\newcommand{\ZusammenfassungEN}{
As part of this Bachelor’s thesis, an automatable system for assembling linear guides was developed for SmarAct GmbH. The aim was to improve the efficiency, reproducibility, and quality stability of the previously manual assembly process. After analyzing the existing process, technical requirements were derived.\\
Based on this analysis, an assembly system was designed, implemented, and tested. Core components include a precise positioning system, automated screw-driving tools, a pressing system for wedge assembly, and a sliding mechanism for preload adjustment. In addition, a suitable control concept for process monitoring and validation was developed. The solution was successfully tested and optimized to enable stable, high-quality serial assembly.\\
This thesis is subject to a confidentiality agreement; therefore, detailed results are not publicly.
}
